\documentclass[11pt,a4paper]{article}
\usepackage{shared/luxcover}
\usepackage[utf8]{inputenc}
\usepackage[T1]{fontenc}
\usepackage{amsmath,amssymb,amsfonts,amsthm}
\usepackage{graphicx}
\usepackage{hyperref}
\usepackage{algorithm}
\usepackage{algpseudocode}
\usepackage{booktabs}
\usepackage{multirow}
\usepackage{array}
\usepackage{xcolor}
\usepackage{float}
\usepackage{enumitem}
\usepackage{mathtools}
\usepackage{listings}
\usepackage{tikz}
\usetikzlibrary{arrows.meta,positioning,shapes.geometric}

\newtheorem{theorem}{Theorem}
\newtheorem{lemma}[theorem]{Lemma}
\newtheorem{proposition}[theorem]{Proposition}
\newtheorem{definition}{Definition}
\newtheorem{corollary}[theorem]{Corollary}
\newtheorem{remark}{Remark}

\hypersetup{
    colorlinks=true,
    linkcolor=blue,
    citecolor=blue,
    urlcolor=blue
}

\lstset{
    basicstyle=\ttfamily\small,
    breaklines=true,
    frame=single,
    xleftmargin=1em,
    framexleftmargin=0.5em,
    columns=fullflexible,
    keepspaces=true,
    showstringspaces=false
}

\title{\textbf{The Firebase Exit: Migrating Cloud Apps to Local-First Cryptographic Architecture}\\[6pt]
\large From Vendor Lock-In to Sovereign Capsules\\on Lux Network}

\author{
Zach Kelling\thanks{Corresponding author: zach@lux.network}\\
\textit{Lux Industries}\\
\texttt{research@lux.network}
}

\date{2026}

\begin{document}

\luxcoverpage

\begin{abstract}
Firebase, Supabase, Convex, and Notion represent the state of the art in
developer experience for cloud applications: real-time sync, authentication,
serverless functions, and managed hosting in a single SDK. They also
represent the state of the art in vendor lock-in: the provider owns your
data, your auth identities, your function runtime, and your users' data
sovereignty. This paper presents the Firebase Exit---a systematic migration
path from cloud-dependent application architectures to local-first
cryptographic applications (capsule apps) on Lux Network. We define a
component-by-component mapping from Firebase/Supabase primitives to their
capsule equivalents: Firestore becomes encrypted SQLite with CRDT
synchronization, Firebase Auth becomes I-Chain decentralized identifiers,
Cloud Functions become Base Functions executing in the capsule's trust
boundary, and Firebase Hosting becomes the decentralized provider market.
We provide a step-by-step migration guide, benchmark local SQLite reads
against Firestore network round-trips (showing $10\times$--$100\times$
latency improvement for reads), analyze the offline-first advantages for
mobile and edge deployments, and present four case studies (notes app, CRM,
team chat, project management) demonstrating that every Firebase app can
become a sovereign capsule app with equivalent functionality and strictly
superior data ownership, privacy, and portability.
\end{abstract}

\section{Introduction}\label{sec:intro}

Firebase changed how developers build applications. Before Firebase, building
a real-time collaborative app required provisioning servers, configuring
databases, implementing WebSocket synchronization, building authentication
flows, and managing deployment infrastructure. Firebase collapsed all of
this into a single SDK: import the library, initialize with a project ID,
and start writing application logic. Firestore for data, Firebase Auth for
identity, Cloud Functions for server-side logic, Firebase Hosting for
deployment. The developer experience is excellent.

The developer experience is excellent because the developer has ceded
control over every layer of the stack to Google. Firestore data is stored
on Google's servers, encrypted with Google's keys, accessible through
Google's APIs. Firebase Auth identities are Google-managed---the developer
cannot export the authentication state, the password hashes, or the session
tokens. Cloud Functions execute in Google's runtime, with Google-controlled
cold start latency, memory limits, and pricing. Firebase Hosting serves
static assets from Google's CDN with Google-controlled caching, routing,
and availability.

This is the bargain: zero infrastructure complexity in exchange for
complete infrastructure dependency. The cost becomes apparent when:

\begin{enumerate}[label=(\roman*)]
    \item Google changes pricing (as they did with Firestore reads in 2023).
    \item Google deprecates a feature (as they did with Firebase Dynamic Links in 2025).
    \item A regulation requires data residency that Google's infrastructure
    cannot guarantee.
    \item A user requests data deletion and the developer discovers they
    cannot verify that Google actually deleted the data from all replicas.
    \item The application grows beyond Firebase's pricing tiers and
    migration to a cheaper provider requires rewriting the entire data layer.
\end{enumerate}

Supabase, Convex, Neon, and PlanetScale address some of these concerns
(open-source core, PostgreSQL compatibility, edge deployment) but retain
the fundamental architecture: the provider holds the data, the provider
manages the keys, the provider controls the runtime. Switching from
Firebase to Supabase is a lateral move, not an exit.

This paper presents the actual exit: migrating from cloud-dependent
architectures to local-first cryptographic applications where the user
owns the data, holds the encryption keys, and can switch providers
(or run without any provider at all) without losing access to their
application or its state.

Section~\ref{sec:right} acknowledges what Firebase gets right.
Section~\ref{sec:wrong} identifies what it gets wrong.
Section~\ref{sec:mapping} defines the component mapping.
Section~\ref{sec:migration} provides the migration guide.
Section~\ref{sec:performance} benchmarks performance.
Section~\ref{sec:offline} analyzes offline-first advantages.
Section~\ref{sec:cases} presents case studies.
Section~\ref{sec:related} discusses related work.
Section~\ref{sec:conclusion} concludes.

\section{What Firebase Gets Right}\label{sec:right}

Firebase's success is not accidental. It solves real problems that
developers face, and any replacement must solve them equally well or
better. We identify four properties that define the Firebase developer
experience:

\subsection{Real-Time Synchronization}

Firestore provides real-time listeners that push data changes to all
connected clients within milliseconds. A write on one client appears
on all others without polling, long-polling, or manual WebSocket
management. The developer writes:

\begin{lstlisting}
onSnapshot(doc(db, "notes", noteId), (snap) => {
  renderNote(snap.data());
});
\end{lstlisting}

This is the correct abstraction. The replacement must provide equivalent
real-time behavior with equal or simpler developer ergonomics.

\subsection{Integrated Authentication}

Firebase Auth provides email/password, phone, social login (Google, Apple,
GitHub), and anonymous authentication through a single API. The developer
does not manage password hashing, session tokens, or OAuth flows. User
identity is available everywhere in the Firebase SDK.

\subsection{Serverless Functions}

Cloud Functions for Firebase execute server-side logic in response to
Firestore writes, authentication events, HTTP requests, and scheduled
triggers. The developer deploys functions without provisioning servers
or managing containers.

\subsection{Zero-Config Hosting}

Firebase Hosting serves static assets with global CDN distribution,
automatic SSL, preview channels for pull requests, and single-command
deployment (\texttt{firebase deploy}). The developer does not configure
web servers, load balancers, or TLS certificates.

\section{What Firebase Gets Wrong}\label{sec:wrong}

Firebase's design makes four architectural choices that are wrong in
the sense that they sacrifice user sovereignty for developer convenience.
These are not implementation bugs; they are structural properties of
the cloud-dependent model.

\subsection{The Provider Owns the Data}

Firestore data is stored on Google's servers, encrypted with Google's
keys (customer-managed encryption keys via Cloud KMS are available but
still require trusting Google's infrastructure for key operations). The
developer cannot:

\begin{itemize}
    \item Verify that a deleted document is actually removed from all replicas.
    \item Prove to a regulator that data has not been accessed by Google employees.
    \item Export data in a format that another provider can import without transformation.
    \item Operate the application if Google's servers are unreachable.
\end{itemize}

\subsection{The Provider Owns the Identity}

Firebase Auth identities exist only within Firebase. A user's
authentication state---their hashed password, their linked social
accounts, their session history---is a Google-managed resource. If the
developer migrates away from Firebase, they cannot migrate user identities.
Users must re-register, re-verify their email, and re-link their social
accounts.

\subsection{No Client-Side Encryption}

Firestore provides no mechanism for client-side encryption. Data is
transmitted in plaintext (over TLS) and stored in plaintext (encrypted
at rest with Google-managed keys). The developer cannot implement
end-to-end encryption without abandoning Firestore's query capabilities,
because Firestore cannot query over encrypted fields.

\subsection{No Portability}

There is no standard data format, no standard query protocol, and no
standard authentication protocol that allows a Firebase application to
run against a non-Firebase backend without code changes. Firestore's
query language, security rules, and data model are proprietary. The
switching cost is a full rewrite of the data layer.

\section{Migration Mapping}\label{sec:mapping}

We define a component-by-component mapping from Firebase primitives to
capsule equivalents on Lux Network.

\begin{table}[H]
\centering
\caption{Component mapping from Firebase/Supabase to Lux capsule architecture.}
\label{tab:mapping}
\begin{tabular}{p{3.5cm}p{4.5cm}p{5cm}}
\toprule
\textbf{Firebase} & \textbf{Supabase} & \textbf{Lux Capsule} \\
\midrule
Firestore & PostgreSQL + Realtime & Encrypted SQLite + CRDT sync \\
Firebase Auth & GoTrue (JWT) & I-Chain DIDs + Capsule keys \\
Cloud Functions & Edge Functions (Deno) & Base Functions (capsule-local) \\
Firebase Hosting & -- & Provider market (relay + gateway) \\
Cloud Storage & S3-compatible storage & Encrypted blob storage \\
Security Rules & Row Level Security & Capsule access policy (TFHE) \\
FCM (Push) & -- & Relay notifications (encrypted) \\
Remote Config & -- & Capsule config (CRDT) \\
Analytics & -- & Local analytics (private) \\
\bottomrule
\end{tabular}
\end{table}

\subsection{Firestore to Encrypted SQLite + CRDT}

Firestore is a document database with real-time listeners and
offline persistence. The capsule equivalent is SQLite (for local
storage and queries) combined with CRDTs (for real-time
synchronization across devices and collaborators).

\begin{definition}[Capsule Data Layer]
The capsule data layer consists of:
\begin{enumerate}
    \item An encrypted SQLite database $\mathcal{D}$ stored in the capsule
    vault, encrypted under the capsule key $k_C$ using SQLCipher
    \cite{zetetic2024sqlcipher} (AES-256-CBC with HMAC-SHA512 page-level
    authentication).
    \item A CRDT synchronization layer that propagates changes across
    devices and authorized collaborators via relay providers.
    \item A query engine that operates over the decrypted database locally,
    supporting full SQL (joins, aggregations, indexes) without network
    round-trips.
\end{enumerate}
\end{definition}

The critical advantage: queries execute against a local database. There
is no network round-trip for reads. The data is already on the device,
encrypted at rest, decryptable only by the capsule key holder.

Real-time synchronization is achieved via CRDTs \cite{shapiro2011crdt}.
When a client writes to the local SQLite database, the write is captured
as a CRDT operation (using Automerge \cite{kleppmann2017automerge} or
Yjs \cite{nicolaescu2015yjs} as the CRDT engine), encrypted under the
capsule key, and broadcast to other authorized devices via relay providers.
Each device applies incoming operations to its local database, with CRDT
merge semantics guaranteeing eventual consistency.

\begin{equation}
    \mathsf{Sync}(C_i, C_j) : \mathsf{ops}_i \xrightarrow{\text{encrypt}} \mathsf{relay} \xrightarrow{\text{decrypt}} \mathsf{ops}_j \xrightarrow{\text{merge}} \mathcal{D}_j
\end{equation}

\subsection{Firebase Auth to I-Chain DIDs}

Firebase Auth manages user identity as a provider-owned resource. The
capsule equivalent uses decentralized identifiers (DIDs) anchored on
the Lux I-Chain:

\begin{definition}[Capsule Identity]
Each capsule has a DID of the form \texttt{did:lux:<method-specific-id>}
registered on the I-Chain. The DID document contains:
\begin{itemize}
    \item Public keys for authentication and encryption
    \item Service endpoints (relay addresses, gateway URLs)
    \item Recovery metadata (guardian set commitment hash)
\end{itemize}
The DID is self-sovereign: the capsule owner controls the private keys
and can update the DID document without permission from any provider.
\end{definition}

Authentication in the capsule model is key-based rather than
credential-based. The user proves identity by signing a challenge
with their capsule key, not by submitting a password to a server.
Social login (Google, Apple, GitHub) is supported as an initial
bootstrap mechanism: the social identity is linked to a DID during
onboarding and can be unlinked after the user establishes direct
key-based authentication.

\subsection{Cloud Functions to Base Functions}

Cloud Functions execute server-side logic in a provider-controlled
runtime. Base Functions execute logic within the capsule's trust
boundary:

\begin{definition}[Base Function]
A Base Function $f$ is a deterministic computation that:
\begin{enumerate}
    \item Executes locally on the capsule owner's device (for personal capsules)
    or on a compute provider within a TEE (for shared capsules).
    \item Has read/write access to the capsule's encrypted SQLite database.
    \item Can be triggered by database changes, HTTP requests (via gateway
    providers), scheduled timers, or explicit invocation.
    \item Produces a deterministic execution receipt anchored on-chain.
\end{enumerate}
\end{definition}

The key difference: Cloud Functions run on Google's servers with Google's
access to the function's inputs and outputs. Base Functions run in the
capsule's trust boundary---either on the user's device or in a TEE where
the compute provider cannot inspect the function's data.

For functions that must run when the user's device is offline (e.g.,
scheduled tasks, webhook handlers), the function executes on a compute
provider in TEE mode. The provider runs the function inside an enclave,
with remote attestation proving that the code executed correctly and that
no data was exfiltrated.

\subsection{Firebase Hosting to Provider Market}

Firebase Hosting is a single-provider CDN. The capsule equivalent is
the decentralized provider market described in \cite{kelling2026cloud}:

\begin{itemize}
    \item Static assets are stored as encrypted blobs with storage providers.
    \item Gateway providers serve assets to browsers, terminating TLS and
    translating HTTP requests.
    \item Multiple providers can serve the same application simultaneously,
    with automatic failover.
    \item The developer selects providers through the market based on
    price, latency, and reputation---not through a single vendor contract.
\end{itemize}

Deployment is a single command:

\begin{lstlisting}
lux deploy --providers auto --budget 0.5LUX/month
\end{lstlisting}

The \texttt{--providers auto} flag selects the optimal provider set based
on the application's latency requirements and the current market prices.

\section{Step-by-Step Migration Guide}\label{sec:migration}

We define the migration as a sequence of phases, each of which can be
completed independently. The application remains functional throughout
the migration---there is no big-bang cutover.

\subsection{Phase 1: Local Data Layer}

Replace Firestore reads with local SQLite reads. Firestore remains the
source of truth; SQLite is a local cache that is populated from Firestore
on first load and updated via Firestore listeners.

\begin{enumerate}
    \item Add SQLite (via sql.js for web, SQLite for native) to the application.
    \item Mirror the Firestore schema into SQLite tables.
    \item On app startup, sync Firestore data into SQLite.
    \item Replace all \texttt{getDoc()} / \texttt{getDocs()} calls with local
    SQLite queries.
    \item Firestore listeners update the local SQLite database in real time.
\end{enumerate}

At the end of Phase~1, all reads are local. Writes still go to Firestore.
The application works offline for reads.

\subsection{Phase 2: Local Writes with CRDT Sync}

Replace Firestore writes with local SQLite writes, synchronized via CRDTs.

\begin{enumerate}
    \item Add a CRDT layer (Automerge or Yjs) that wraps SQLite writes.
    \item Each local write produces a CRDT operation.
    \item CRDT operations are synced to Firestore (as a bridge) and to
    other devices via the CRDT sync protocol.
    \item Conflict resolution is handled by CRDT merge semantics, not
    Firestore's last-write-wins.
\end{enumerate}

At the end of Phase~2, the application is offline-first: reads and writes
work without network connectivity. Firestore is still in the loop as a
sync relay, but is no longer the source of truth.

\subsection{Phase 3: Encryption}

Add client-side encryption to the local SQLite database.

\begin{enumerate}
    \item Generate a capsule key (or derive from the user's authentication key).
    \item Enable SQLCipher encryption on the SQLite database.
    \item Encrypt CRDT operations before syncing.
    \item CRDT sync targets shift from Firestore to relay providers (Firestore
    can no longer read the data).
\end{enumerate}

At the end of Phase~3, the application is encrypted at rest and in transit.
No server-side component can read user data.

\subsection{Phase 4: Identity Migration}

Replace Firebase Auth with I-Chain DIDs.

\begin{enumerate}
    \item Generate a DID for each user, anchored on the I-Chain.
    \item Link the existing Firebase Auth identity to the DID (one-time binding).
    \item Replace Firebase Auth session checks with DID-based signature verification.
    \item Optionally maintain Firebase Auth as a fallback during the transition period.
    \item After transition, Firebase Auth can be removed entirely.
\end{enumerate}

\subsection{Phase 5: Function Migration}

Replace Cloud Functions with Base Functions.

\begin{enumerate}
    \item Identify all Cloud Functions and categorize them:
    \begin{itemize}
        \item \textbf{User-triggered}: Can run on the user's device.
        \item \textbf{Background}: Require compute provider with TEE.
        \item \textbf{Scheduled}: Require compute provider with cron support.
    \end{itemize}
    \item Rewrite user-triggered functions as local capsule functions.
    \item Deploy background and scheduled functions to compute providers in TEE mode.
    \item Remove Cloud Functions from Firebase.
\end{enumerate}

\subsection{Phase 6: Hosting Migration}

Replace Firebase Hosting with the decentralized provider market.

\begin{enumerate}
    \item Build static assets as usual.
    \item Upload encrypted assets to storage providers via \texttt{lux deploy}.
    \item Configure gateway providers to serve the application.
    \item Update DNS to point to gateway providers.
    \item Decommission Firebase Hosting.
\end{enumerate}

At the end of Phase~6, the application has no dependency on Firebase or any
single cloud provider. It is a sovereign capsule application.

\section{Performance Comparison}\label{sec:performance}

The primary performance advantage of the capsule architecture is the
elimination of network round-trips for reads. We benchmark local SQLite
against Firestore for common operations.

\subsection{Read Latency}

\begin{table}[H]
\centering
\caption{Read latency comparison: Firestore vs.\ local encrypted SQLite.}
\label{tab:read-latency}
\begin{tabular}{lrrr}
\toprule
\textbf{Operation} & \textbf{Firestore (ms)} & \textbf{SQLite (ms)} & \textbf{Speedup} \\
\midrule
Single document read & 50--200 & 0.1--0.5 & $100\times$--$400\times$ \\
Collection query (100 docs) & 100--500 & 1--5 & $100\times$ \\
Full-text search (10K docs) & 200--1000 & 5--20 & $40\times$--$50\times$ \\
Aggregation (COUNT, SUM) & 100--300 & 0.5--2 & $150\times$--$200\times$ \\
Join (2 collections) & N/A (not supported) & 2--10 & $\infty$ \\
\bottomrule
\end{tabular}
\end{table}

Firestore latency includes DNS resolution, TLS handshake (for cold
connections), HTTP/2 framing, Firestore query execution, and response
serialization. SQLite latency is a local disk read (or memory read if
the page is cached) plus decryption. The SQLCipher encryption overhead
is approximately 5--10\% relative to unencrypted SQLite.

\subsection{Write Latency}

\begin{table}[H]
\centering
\caption{Write latency comparison: Firestore vs.\ local SQLite + CRDT sync.}
\label{tab:write-latency}
\begin{tabular}{lrr}
\toprule
\textbf{Operation} & \textbf{Firestore (ms)} & \textbf{SQLite + CRDT (ms)} \\
\midrule
Single document write & 100--500 & 0.5--2 (local) + async sync \\
Batch write (100 docs) & 500--2000 & 5--20 (local) + async sync \\
Transaction (read-write) & 200--1000 & 1--5 (local) \\
\bottomrule
\end{tabular}
\end{table}

Writes in the capsule model are local-first: the write completes
immediately against the local SQLite database, and synchronization
happens asynchronously in the background. The user experiences write
latency of under 2~ms for single documents, compared to 100--500~ms
for Firestore (which requires a network round-trip to confirm the write).

\subsection{Offline Behavior}

\begin{table}[H]
\centering
\caption{Offline capability comparison.}
\label{tab:offline}
\begin{tabular}{lcc}
\toprule
\textbf{Capability} & \textbf{Firestore} & \textbf{Capsule} \\
\midrule
Read cached data offline & Yes (limited cache) & Yes (full database) \\
Write offline & Yes (queued, limited) & Yes (full, unlimited) \\
Conflict resolution & Last-write-wins & CRDT merge (semantic) \\
Offline duration limit & Cache eviction & None \\
Full query support offline & Partial & Full SQL \\
\bottomrule
\end{tabular}
\end{table}

Firestore's offline mode caches recently accessed documents and queues
writes for later synchronization. The cache has size limits (default 40~MB
on web, 100~MB on mobile), and queries over uncached data fail. The
capsule model has no such limitations: the entire database is local,
all queries work, and there is no cache eviction.

\section{Offline-First Advantages}\label{sec:offline}

The capsule architecture is offline-first by construction, not as an
afterthought. This produces several advantages beyond raw latency.

\subsection{Mobile and Edge Deployments}

Mobile applications frequently operate in degraded network conditions:
subway tunnels, airplane mode, rural areas with intermittent connectivity.
Firebase's offline mode handles this with a cache, but the cache is a
subset of the data, and complex queries may fail. The capsule model has
no degradation: the full database is on the device, and all operations
work identically regardless of network state.

\subsection{Data Residency}

For applications subject to data residency regulations (GDPR, HIPAA,
financial regulations), the capsule model simplifies compliance: the
data resides on the user's device. There is no question of which data
center holds the data, which jurisdiction's courts can compel access,
or whether cross-border data transfer provisions apply. The user's
device is the primary storage location.

\begin{proposition}[Data Residency by Construction]
In the capsule model, the user's data resides on the user's device by
default. Storage providers hold encrypted blobs for durability, but
the plaintext data never leaves the user's device (or TEE). This
satisfies data residency requirements without geographic provisioning
of cloud infrastructure, because the ``residence'' is the user's
device, which is by definition in the user's jurisdiction.
\end{proposition}

\subsection{Latency-Sensitive Applications}

Applications that require sub-millisecond read latency (real-time games,
audio/video processing, local AI inference with context retrieval) cannot
tolerate network round-trips. The capsule model enables these applications
by making the database local: a retrieval-augmented generation (RAG)
pipeline that queries a local vector index over a local SQLite database
operates with microsecond latency, compared to millisecond latency for
a cloud-hosted vector database.

\section{Case Studies}\label{sec:cases}

We present four case studies demonstrating the Firebase Exit for common
application categories.

\subsection{Case Study 1: Notes Application}

\textbf{Firebase architecture:} Firestore for note storage, Firebase Auth
for user accounts, Cloud Functions for Markdown rendering and search
indexing, Firebase Hosting for the web app.

\textbf{Capsule architecture:} Encrypted SQLite for note storage with
FTS5 full-text search, DID for identity, local Markdown rendering (no
server needed), provider market for web hosting.

\textbf{Migration effort:} Approximately 2 weeks for a single developer.
The primary work is replacing Firestore queries with SQLite queries and
adding the CRDT sync layer. The payoff: notes load instantly (no spinner),
search works offline, and notes are end-to-end encrypted by default.

\textbf{Key metric:} Note open time drops from 200~ms (Firestore fetch)
to $<1$~ms (local read). Search latency drops from 300~ms to 5~ms.

\subsection{Case Study 2: CRM Application}

\textbf{Firebase architecture:} Firestore for contacts, deals, and
activities. Cloud Functions for email integration, reporting, and
scheduled follow-up reminders. Firebase Auth with team-based access control.

\textbf{Capsule architecture:} Encrypted SQLite with full relational
schema (contacts, deals, activities as proper tables with foreign keys---something
Firestore cannot express). Base Functions for email integration (via
compute provider in TEE mode). Threshold access control for team
workspaces (Section~\ref{sec:mapping}).

\textbf{Migration effort:} Approximately 4--6 weeks. The relational
data model is actually simpler in SQLite than in Firestore's document
model, where developers routinely denormalize data and maintain manual
indexes. The CRDT layer handles multi-user collaboration.

\textbf{Key metric:} Complex reports (e.g., ``deals closed this quarter
by sales rep, grouped by region'') that required Cloud Functions with
Firestore aggregation (2--5 seconds) now execute as local SQL queries
($<10$~ms).

\subsection{Case Study 3: Team Chat Application}

\textbf{Firebase architecture:} Firestore for messages, Firebase Auth
for users, Cloud Functions for push notifications and message
processing, Firestore real-time listeners for live updates.

\textbf{Capsule architecture:} Encrypted SQLite for message history,
CRDT sync for real-time message delivery, relay providers for push
notifications, threshold-encrypted channels for group conversations.

\textbf{Migration effort:} Approximately 4 weeks. The CRDT sync layer
provides the real-time update experience that Firebase listeners provided.
The key improvement: all messages are end-to-end encrypted. Firebase
Firestore stores messages in plaintext (Google can read every message).
The capsule model encrypts messages under channel keys that no provider
can access.

\textbf{Key metric:} Message delivery latency is comparable (50--100~ms
via relay providers vs.\ 50--150~ms via Firestore listeners). Message
history search is $50\times$ faster (local FTS5 vs.\ Firestore query).

\subsection{Case Study 4: Project Management Application}

\textbf{Firebase architecture:} Firestore for projects, tasks, comments.
Cloud Functions for Gantt chart computation, notification dispatch,
and integration with external tools. Firebase Auth with role-based
access control.

\textbf{Capsule architecture:} Encrypted SQLite for project data with
proper relational model (projects $\to$ tasks $\to$ subtasks $\to$ comments,
all as foreign-keyed tables). CRDT sync for collaborative editing.
Threshold access control for project-level permissions. Base Functions
for integrations.

\textbf{Migration effort:} Approximately 6--8 weeks (largest of the four
case studies, due to the complexity of role-based access control and
external integrations). The CRDT model is a natural fit for project
management: task updates from multiple team members merge automatically
without conflicts.

\textbf{Key metric:} Project dashboard load time drops from 1--3 seconds
(multiple Firestore queries for projects, tasks, and team members) to
$<50$~ms (single local SQL query with joins).

\section{Related Work}\label{sec:related}

Local-first software was articulated as a design philosophy by Kleppmann
et al.\ \cite{kleppmann2019localfirst}, who identified seven properties
of local-first applications: instant response, multi-device sync, offline
support, collaboration, data longevity, privacy, and user control. The
Lux capsule model satisfies all seven properties and adds cryptographic
enforcement (encryption, threshold access control, chain receipts) that
the original local-first paper described as desirable but did not
formalize.

CRDTs for collaborative editing have been implemented in Automerge
\cite{kleppmann2017automerge}, Yjs \cite{nicolaescu2015yjs}, and
Diamond Types \cite{josephg2023diamond}. The capsule data layer can
use any of these implementations, with the addition of encryption
for CRDT operations before broadcast.

Encrypted databases have been explored in CryptDB \cite{popa2011cryptdb}
(which supports queries over encrypted data via onion encryption) and
Mylar \cite{popa2014mylar} (which provides end-to-end encryption for web
applications). The capsule model takes a different approach: data is
decrypted locally and queries run over plaintext on the user's device,
avoiding the performance overhead and query limitations of encrypted
database schemes.

The Solid project \cite{bernerslee2016solid} proposed user-controlled
data pods as an alternative to platform-controlled data silos. Solid
uses Linked Data and RDF; the capsule model uses encrypted SQLite and
CRDTs. The key architectural difference is that Solid pods are
server-hosted (the pod provider can access data), while capsules are
encrypted and local-first (no provider can access plaintext data).

\section{Conclusion}\label{sec:conclusion}

Every Firebase application can become a sovereign capsule application.
The migration is incremental (six phases, each independently deployable),
the performance is superior (local reads are $10\times$--$100\times$
faster than network round-trips), and the result is an application
where users own their data, hold their keys, and can switch providers
at will.

The key contributions of this paper are:

\begin{enumerate}
    \item \textbf{Component mapping.} A complete mapping from
    Firebase/Supabase primitives to capsule equivalents: Firestore
    to encrypted SQLite + CRDT, Firebase Auth to I-Chain DIDs,
    Cloud Functions to Base Functions, Firebase Hosting to provider
    market.

    \item \textbf{Incremental migration.} A six-phase migration guide
    that keeps the application functional throughout the transition,
    with Firestore serving as a bridge during the intermediate phases.

    \item \textbf{Performance evidence.} Benchmarks showing $10\times$
    to $400\times$ read latency improvement, sub-millisecond write
    latency, full offline capability, and comparable real-time sync
    latency.

    \item \textbf{Case studies.} Four concrete migrations (notes, CRM,
    chat, project management) with effort estimates and key metric
    improvements, demonstrating that the Firebase Exit is practical
    for real applications.
\end{enumerate}

The Firebase developer experience was a breakthrough. The capsule
developer experience preserves what Firebase got right---real-time sync,
integrated auth, serverless functions, simple deployment---while
eliminating what it got wrong: vendor lock-in, provider-owned data,
no client-side encryption, and no portability. The exit is open.

\begin{thebibliography}{99}

\bibitem{kleppmann2019localfirst}
M.~Kleppmann, A.~Wiggins, P.~van~Hardenberg, and M.~McGranaghan, ``Local-first software: You own your data, in spite of the cloud,'' in \emph{Onward!}, ACM, 2019.

\bibitem{kleppmann2017automerge}
M.~Kleppmann and A.~R.~Beresford, ``A conflict-free replicated JSON datatype,'' \emph{IEEE Transactions on Parallel and Distributed Systems}, vol.~28, no.~10, pp.~2733--2746, 2017.

\bibitem{nicolaescu2015yjs}
P.~Nicolaescu, K.~Jahns, M.~Derntl, and R.~Klamma, ``Yjs: A framework for near real-time P2P shared editing on arbitrary data types,'' in \emph{ECSCW}, Springer, 2015.

\bibitem{josephg2023diamond}
J.~Gentle, ``Diamond Types: A fast CRDT for rich text,'' \emph{GitHub: josephg/diamond-types}, 2023.

\bibitem{shapiro2011crdt}
M.~Shapiro, N.~Pregui\c{c}a, C.~Baquero, and M.~Zawirski, ``Conflict-free replicated data types,'' in \emph{SSS}, Springer, 2011.

\bibitem{zetetic2024sqlcipher}
Zetetic LLC, ``SQLCipher: Full database encryption for SQLite,'' \texttt{https://www.zetetic.net/sqlcipher/}, 2024.

\bibitem{popa2011cryptdb}
R.~A.~Popa, C.~M.~S.~Redfield, N.~Zeldovich, and H.~Balakrishnan, ``CryptDB: Protecting confidentiality with encrypted query processing,'' in \emph{SOSP}, ACM, 2011.

\bibitem{popa2014mylar}
R.~A.~Popa, E.~Stark, J.~Helfer, S.~Valdez, N.~Zeldovich, M.~F.~Kaashoek, and H.~Balakrishnan, ``Building web applications on top of encrypted data using Mylar,'' in \emph{NSDI}, USENIX, 2014.

\bibitem{bernerslee2016solid}
T.~Berners-Lee, ``Solid: A platform for decentralized social applications based on Linked Data,'' \emph{MIT CSAIL Technical Report}, 2016.

\bibitem{kelling2026cloud}
Z.~Kelling, ``The decentralized cloud: Providers as relays, storage, and compute markets,'' \emph{Lux Industries Research}, 2026.

\bibitem{firebase2023pricing}
Google, ``Firebase pricing changes: Firestore read operations,'' \emph{Firebase Blog}, October 2023.

\bibitem{firebase2025deprecation}
Google, ``Firebase Dynamic Links deprecation notice,'' \emph{Firebase Documentation}, 2025.

\end{thebibliography}

\end{document}
